% !PRINTED 291
\weberpage{291}
%<<CONT
thus the characteristic is equal to half the excess of the first number over the second.

To this theorem we can give the following expression, in which it appears entirely independent of relations of magnitude and depends only on the relations of position of the curves and of their points of intersection.

If we distinguish the points of intersection of $\varphi$, $\psi$, according as they are exit points $A(\varphi,\psi)$ or entry points $E(\varphi,\psi)$, by the signs $\alpha$ and $\varepsilon$, and likewise, according as they are exterior or interior points with respect to the curve $f$, by $\alpha$ and $\varepsilon$, and then give the points $\alpha,\varepsilon$ and $\varepsilon,\alpha$ the \gesperrt{character $+1$, and the points $\alpha,\alpha$ and $\varepsilon,\varepsilon$ the character $-1$}\textsuperscript{1)}, \gesperrt{then the characteristic of the system of functions $(f,\varphi,\psi)$ is equal to half the sum of the characters of all the points of intersection of $\varphi,\psi$.}

The theorems and definitions that we have given here remain valid unchanged even when the curves considered consist of several closed circuits; double points may also be present on the curves; it is only necessary that, for every direction of advance on one of the curves, it be completely determined whether it is to be regarded as positive or as negative, i.\,e. every piece of a curve $f$ or $\varphi$ or $\psi$ must separate portions of surface from one another in which the functions $f$ or $\varphi$ or $\psi$ have opposite signs.

Only those cases are to be excluded in which one of the curves passes through a double point of the other, or in which the curves touch one another, or in which the three curves $f$, $\varphi$, $\psi$ pass through \gesperrt{one} point; for in these cases the determination of a point as an entry point or exit point would become doubtful.

If the curves extend to infinity, then, in order to make our theorems applicable, one must close them off by arbitrarily added pieces of curve, as we shall see in the next section.

\noindent\rule{4cm}{0.4pt}

\noindent\textsuperscript{1)} \gesperrt{Kronecker} understands by the character of a point something different.
